%% toc: P2 15 15.5,15.6,15.7
\section{Proofs for the least-model link and Cook--Levin}
\label{sec:annex-p02-w02b-p02-015-links}

Clause lists, \(A_S\), bodies, heads, closed sets and \(v_X\) are as in
Convention~\ref{def:p02-w02-p02-019-clause-input} and
Definition~\ref{def:p02-w02-p02-019-horn-parts}; \(P_S\), \(T_{P_S}\) and
\(M_{P_S}\) are as in Definition~\ref{def:p02-w02b-p02-015-horn-program}.
Machines are read as in
Convention~\ref{def:p02-w02b-p02-015-machine-reading}.

\begin{annexitem}{annex:p02-w02b-p02-015-horn-lfp}{thm:p02-w02b-p02-015-horn-lfp}{Forward closure is the least model of the definite program}
\begin{proof}
Write \(T=T_{P_S}\) and \(F=\operatorname{Fwd}(S)\).

\emph{Clause 1.}  By Definition~\ref{def:p09-definite-program},
\(T(I)\) consists of the atoms \(H\theta\) for program clauses
\(H\leftarrow B_1,\ldots,B_k\) of \(P_S\) and ground substitutions \(\theta\)
with every \(B_i\theta\in I\).  The program clauses contain no variable, so
\(H\theta=H\) and \(B_i\theta=B_i\) for every \(\theta\).  By
Definition~\ref{def:p02-w02b-p02-015-horn-program} the program clauses are
exactly the clauses \(\operatorname{hd}(C)\leftarrow q_1,\ldots,q_k\) with
\(C\) a definite clause of \(S\) and \(\{q_1,\ldots,q_k\}=B(C)\).  Hence
\(T(I)\) is the set of heads of definite clauses \(C\) of \(S\) with
\(B(C)\subseteq I\).  If \(I\subseteq J\) and \(B(C)\subseteq I\), then
\(B(C)\subseteq J\); so \(T(I)\subseteq T(J)\).

\emph{Clause 2.}  By Definition~\ref{def:p02-w02-p02-019-horn-parts}, \(X\)
is closed under \(S\) iff \(\operatorname{hd}(C)\in X\) for every definite
clause \(C\) of \(S\) with \(B(C)\subseteq X\).  By clause~1 this says
exactly that every member of \(T(X)\) lies in \(X\).

\emph{Clause 3.}  Put \(U_k=T^k(\varnothing)\) and
\(U=\bigcup_{k<\omega}U_k\).  First, \(U_k\subseteq U_{k+1}\) by induction
on \(k\): \(U_0=\varnothing\subseteq U_1\); and if \(U_k\subseteq U_{k+1}\),
monotonicity gives \(U_{k+1}=T(U_k)\subseteq T(U_{k+1})=U_{k+2}\).

Second, \(U_k\subseteq X\) for every \(k\) and every \(X\subseteq A_S\) with
\(T(X)\subseteq X\), again by induction: \(U_0=\varnothing\subseteq X\), and
if \(U_k\subseteq X\), then \(U_{k+1}=T(U_k)\subseteq T(X)\subseteq X\).
Hence \(U\subseteq X\) for every such \(X\).

Third, \(T(U)=U\).  Let \(a\in T(U)\).  By clause~1, \(a=\operatorname{hd}(C)\)
for a definite clause \(C\) with \(B(C)\subseteq U\).  The set \(B(C)\) is
finite, each of its members lies in some \(U_k\), and the \(U_k\) increase,
so \(B(C)\subseteq U_j\) for the largest such index \(j\) (or for \(j=0\) if
\(B(C)=\varnothing\)).  Then \(a\in T(U_j)=U_{j+1}\subseteq U\).  Conversely,
let \(a\in U\), say \(a\in U_k\).  Then \(k\geq1\), because
\(U_0=\varnothing\), and \(a\in U_k=T(U_{k-1})\subseteq T(U)\) by
monotonicity, since \(U_{k-1}\subseteq U\).  Every \(U_k\) is a subset of
\(A_S\), because every head of a clause of \(S\) is a variable of \(S\); so
\(U\subseteq A_S\).

Now \(T(U)\subseteq U\) and clause~2 give that \(U\) is closed under \(S\),
so \(F\subseteq U\) because \(F\) is the least closed subset of \(A_S\)
(Definition~\ref{def:p02-w02-p02-019-forward-closure}).  Also \(F\) is closed,
so \(T(F)\subseteq F\) by clause~2, and the second step gives
\(U\subseteq F\).  So \(F=U\), and \(T(F)=F\) by the third step.  The set
\(M_{P_S}\) is by Definition~\ref{def:p09-definite-program} the union
\(\bigcup_{k<\omega}T^k(\varnothing)=U\); it is the least fixed point of
\(T\), because \(T(U)=U\) and every fixed point \(X\) satisfies
\(T(X)\subseteq X\) and hence contains \(U\).  This proves
\(F=U=M_{P_S}\), and the second step gives \(F\subseteq X\) whenever
\(T(X)\subseteq X\).

For the stage bound, let \(j\) be the least index with \(U_{j+1}=U_j\).  Such
an index exists and satisfies \(j\leq|A_S|\): if \(U_0\subsetneq
U_1\subsetneq\cdots\subsetneq U_i\), then \(|U_i|\geq i\), and \(|U_i|\leq
|A_S|\), so a strictly increasing run has at most \(|A_S|\) steps.  If
\(U_{j+1}=U_j\), then \(U_{j+2}=T(U_{j+1})=T(U_j)=U_{j+1}\), and by induction
\(U_k=U_j\) for all \(k\geq j\).  Hence \(U=U_j\), and \(U_k=U=F\) for every
\(k\geq|A_S|\geq j\).

\emph{Clause 4.}  Let \(I\subseteq A_S\).  The valuation \(v_I\) satisfies a
definite clause \(C\) iff \(v_I(\operatorname{hd}(C))=\mathbf T\) or some
\(q\in B(C)\) has \(v_I(q)=\mathbf F\)
(Definition~\ref{def:p02-w02-p02-019-horn-parts}, item~1), that is, iff
\(B(C)\subseteq I\) implies \(\operatorname{hd}(C)\in I\).  So \(v_I\)
satisfies every definite clause of \(S\) iff \(I\) is closed under \(S\),
iff \(T(I)\subseteq I\) by clause~2.  A Herbrand interpretation
\(I\subseteq\mathit{HB}_{P_S}=A_S\) is a model of \(P_S\) iff each program
clause \(r\leftarrow q_1,\ldots,q_k\) with all \(q_i\in I\) has \(r\in I\),
which is the same condition.  By clause~3, \(M_{P_S}=F\) satisfies
\(T(F)\subseteq F\), so it is a Herbrand model, and it is contained in every
\(I\) with \(T(I)\subseteq I\), that is, in every Herbrand model.  So it is
the least Herbrand model.

For the last sentence, let \(a\in A_S\).  If \(a\in M_{P_S}=F\) and \(v\)
satisfies every definite clause of \(S\), put
\(X_v=\{p\in A_S:v(p)=\mathbf T\}\).  The valuations \(v\) and \(v_{X_v}\)
agree on \(A_S\), so \(v_{X_v}\) satisfies the same clauses of \(S\), and by
the first part of this clause \(T(X_v)\subseteq X_v\).  By clause~3,
\(F\subseteq X_v\), so \(v(a)=\mathbf T\).  If \(a\notin M_{P_S}\), then
\(v_{M_{P_S}}\) satisfies every definite clause of \(S\), by the first part
of this clause and \(T(M_{P_S})=M_{P_S}\), and it makes \(a\) false.

\emph{Clause 5.}  By clause~3, \(M_{P_S}=F\).  Clause~4 of
Theorem~\ref{thm:p02-w02-p02-019-horn-closure} states that \(S\) is
unsatisfiable iff \(S\) contains a negative clause \(C\) with
\(B(C)\subseteq F\), and that, when \(S\) is satisfiable, \(v_F\) satisfies
\(S\) and every valuation satisfying \(S\) makes true every atom that \(v_F\)
makes true.  Substituting \(M_{P_S}\) for \(F\) gives clause~5.  A negative
clause \(C\) with \(B(C)=\{q_1,\ldots,q_k\}\) is falsified by \(v_I\) iff
every \(q_i\in I\), that is, iff \(I\) violates the goal
\(\leftarrow q_1,\ldots,q_k\); this is the reading in terms of goals.
\end{proof}
\end{annexitem}

\begin{annexitem}{annex:p02-w02b-p02-015-bridge}{thm:p02-w02b-p02-015-bridge}{Part~2 and Part~4 conventions for maps, reductions and \(\mathsf{NP}\)}
\begin{proof}
\emph{Marked cells.}  The simulations below extend the tape alphabet of a
machine by a marked copy \(\hat a\) of each symbol \(a\), the blank included.
A marked copy is written only in cell \(0\) of a tape, before the
simulation proper starts.  A head that must return to cell \(0\) moves left
until it scans a marked symbol, which takes at most as many steps as its
distance from cell \(0\).  During a simulation every scanned \(\hat a\) is
treated as \(a\), and every symbol written is unmarked.

\emph{Clause 1.}  A machine with a read-only input tape is a machine whose
input tape is written only with the scanned symbol, so the implication from
read-only to writable is immediate, with the same number of steps.
Conversely, let \(M\) have tapes \(1,\ldots,\kappa\), let its output tape be
tape \(o\neq1\), and let it halt within \(p(n)\) steps on inputs of length
\(n\).  The machine \(M'\) has a read-only tape~\(1\) and work tapes
\(2,\ldots,\kappa+1\).  On input \(x\) of length \(n\) it first copies
\(x\) into cells \(0,\ldots,n-1\) of tape~\(2\), moving the heads of tapes
\(1\) and \(2\) right together and writing the symbol in cell \(0\) of
tape~\(2\) in marked form; if \(n=0\), it writes \(\hat\sqcup\) in that cell.
This takes at most \(n+1\) steps.  It then returns the head of tape~\(2\) to
cell \(0\) in at most \(n+1\) steps.  Tapes \(3,\ldots,\kappa+1\) are blank
with their heads on cell \(0\).  Now \(M'\) simulates \(M\) one step per
step, with tape \(j+1\) of \(M'\) in the role of tape \(j\) of \(M\).  The
configuration of \(M'\) restricted to tapes \(2,\ldots,\kappa+1\) is, after
the copying phase and after each simulated step, the configuration of \(M\)
after the same number of its own steps, read with marked symbols as
unmarked; this holds at the start by construction and is preserved because
\(M'\) applies \(\delta\) of \(M\) to the unmarked readings of the scanned
symbols.  So \(M'\) halts in the state that \(M\) halts in, within
\(2n+2+p(n)\) steps, and its tape \(o+1\) then holds what tape \(o\) of \(M\)
holds.  Since \(o\neq1\), tape \(o+1\) is not tape~\(2\), carries no marked
symbol, and holds the output of \(M\) in the required format.  The bound
\(2n+2+p(n)\) is a polynomial.  For deciding a language the output tape plays no role.

Since both reduction notions require a polynomial-time computable map
\(f\) with \(x\in A\) iff \(f(x)\in B\), and the two computability notions
agree, the two reduction relations are the same.  The class \(\mathsf P\) of
Definition~\ref{def:p04-resource-classes} uses
\(\bigcup_k\mathsf{DTIME}(n^k)\), whose machines halt within \(c\,n^k\) steps
for all \(n\geq n_0\); the finitely many shorter inputs take at most some
fixed number \(d\) of steps, so such a machine halts within \(c\,n^k+d\)
steps on every input.  Conversely a polynomial bound \(p(n)\leq
e\,(n+1)^k\leq e\,2^kn^k\) for \(n\geq1\) gives membership in
\(\mathsf{DTIME}(n^k)\).  So the two definitions of \(\mathsf P\) coincide.

\emph{Clause 2.}  Let \(N\) be a nondeterministic machine with read-only
input that decides \(L\) and whose running time is \(O(n^k)\): there are
\(c,n_0\) such that every branch on an input of length \(n\geq n_0\) has at
most \(c\,n^k\) steps.  For each of the finitely many inputs \(x\) of length
below \(n_0\), every branch is finite because \(N\) is a decider, and the
lengths of the branches on \(x\) are bounded: otherwise the tree of finite
initial segments of branches, in which every node has at most
\(\max_{s,\vec a}|\Delta(s,\vec a)|\) children, would have infinitely many
nodes, and choosing at each node a child with infinitely many descendants
would produce an infinite branch.  Let \(d\) bound these lengths for all
inputs of length below \(n_0\), and put \(t(n)=c\,n^k+d\).  Every branch on
every input of length \(n\) has at most \(t(n)\) steps.

For every nonhalting state \(s\) and tuple \(\vec a\) list
\(\Delta(s,\vec a)\) in a fixed order as \(O_0,\ldots,O_{r-1}\), where
\(r=r(s,\vec a)\).  Let \(b\geq1\) bound every \(r(s,\vec a)\), and let
\(\beta\geq1\) be least with \(2^\beta\geq b\).  A bit string \(y\) is read
as a sequence of blocks of \(\beta\) bits, the block of index \(i\) occupying
positions \(i\beta+1,\ldots,(i+1)\beta\); positions beyond the end of \(y\)
read as \(0\).  For a bit string \(y\), let \(\rho_y\) be the sequence of
configurations that starts with the initial configuration of \(N\) on \(x\),
stops at a configuration with a halting state or with \(r=0\), and otherwise
passes from its \(i\)th configuration, with state \(s\) and scanned tuple
\(\vec a\), by the triple \(O_{j\bmod r}\), where \(j\) is the value in
binary of block \(i\) of \(y\).  Every step of \(\rho_y\) uses a triple of
\(\Delta\), so \(\rho_y\) is a branch of \(N\) on \(x\), and it has at most
\(t(n)\) steps.  Define \(R(x,y)\) to hold iff \(\rho_y\) ends in
\(q_{\mathrm a}\), and put \(q(n)=\beta\,t(n)\), a polynomial.

\emph{Equivalence.}  If \(R(x,y)\) for some \(y\), then \(\rho_y\) is an
accepting branch, so \(N\) accepts \(x\) and \(x\in L\).  If \(x\in L\), some
branch accepts; let it have \(T'\leq t(n)\) steps and use at step \(i\) the
triple \(O_{j_i}\) with \(j_i<r\leq b\leq2^\beta\).  Let \(y\) be the
concatenation of the \(\beta\)-bit binary numerals of
\(j_0,\ldots,j_{T'-1}\).  Then \(j_i\bmod r=j_i\), so \(\rho_y\) is that
branch, \(R(x,y)\) holds, and \(\operatorname{len}(y)=\beta T'\leq q(n)\).
So \(x\in L\) iff some \(y\) of length at most \(q(n)\) satisfies
\(R(x,y)\).

\emph{\(R\) is polynomial time.}  A deterministic machine \(M_R\) with
tapes \(1,\ldots,\kappa_N+3\), where \(\kappa_N\) is the number of tapes of
\(N\), decides \(\{\langle x,y\rangle:R(x,y)\}\) as follows, on an input
\(w\) of length \(\ell\).  (i) It reads \(w\) in consecutive pairs of bits.
While the pair is \(00\) or \(11\) it writes the repeated bit on tape~\(2\)
and moves right; at the first pair \(01\) it copies every remaining bit of
\(w\) onto tape~\(3\).  The first symbol written on each of tapes \(2\) and
\(3\) is marked, and a marked blank is written in cell \(0\) of a tape that
receives no bit.  If a pair \(10\) occurs before the first \(01\), or the
input ends before a pair \(01\) is read, \(M_R\) rejects, since then
\(w\) is no pair \(\langle x,y\rangle\)
(Convention~\ref{def:p02-w02-p02-018-verifier}).  This phase takes at most
\(\ell+1\) steps.  (ii) It returns the heads of tapes \(2\) and \(3\) to cell
\(0\), in at most \(\ell\) steps.  Tape~\(2\) now holds \(x\), tape~\(3\)
holds \(y\), and \(n\leq\ell\).  (iii) It simulates \(\rho_y\), with tape~\(2\)
as the read-only input tape of \(N\) and tapes \(4,\ldots,\kappa_N+3\) as the
work tapes of \(N\); tapes \(2\) and \(3\) are never written in this phase.
For each simulated step it moves the head of tape~\(3\) right across the next
\(\beta\) cells, recording the block value in its finite control (a blank
reads as \(0\)), and then applies \(O_{j\bmod r}\) to the simulated tapes in
one further step; \(r\), the order of \(\Delta(s,\vec a)\) and \(j\bmod r\)
are finite data of \(N\) stored in the control.  It accepts when the
simulated state is \(q_{\mathrm a}\), and rejects at \(q_{\mathrm r}\) or
when \(r=0\).  Since \(\rho_y\) has at most \(t(n)\) steps, \(M_R\) halts
within
\[
 (\ell+1)+\ell+(\beta+1)\,t(n)+1\leq 2\ell+2+(\beta+1)\,(c\,\ell^k+d)
\]
steps, a polynomial in \(\ell\).  So \(R\) is a polynomial-time predicate and
\(L\in\mathsf{NP}\) in the sense of
Convention~\ref{def:p02-complexity-standing}.

\emph{Clause 3.}  Let \(M_R\) decide \(\{\langle x,y\rangle:R(x,y)\}\) with
tapes \(1,\ldots,\kappa\) within \(p(\ell)\) steps, where, after enlarging
\(p\), \(p(\ell)=e\,(\ell+1)^d\) with \(e,d\geq1\), so that \(p\) is
nondecreasing.  The nondeterministic machine \(N\) has a read-only tape~\(1\)
and work tapes \(2,\ldots,\kappa+2\).  On input \(x\) of length \(n\):

\emph{Phase 1.}  It writes \(\hat{\triangleright}\) in cell \(0\) of
tape~\(3\) and moves that head right, where \(\triangleright\) is a new
symbol.  Then, for each bit \(x_i\) under the input head, it writes \(x_i\)
on tape~\(2\) and moves right, writes \(x_i\) again and moves right while the
input head moves right, and writes \(1\) on tape~\(3\) and moves right.  On
reaching the blank after \(x\), it writes \(0\) and then \(1\) on tape~\(2\).
The first symbol written on tape~\(2\) is marked.  Tape~\(2\) now holds
\(x_1x_1\cdots x_nx_n01\), the first part of \(\langle x,y\rangle\), with
its head on the next cell, and tape~\(3\) holds \(\hat\triangleright\)
followed by \(n\) cells containing \(1\).  This takes at most \(2n+3\) steps.

\emph{Phase 2.}  It moves the head of tape~\(3\) one cell left, onto cell
\(n\).  Then it repeats: if the head of tape~\(3\) scans
\(\hat\triangleright\), it goes to Phase~3; otherwise it chooses one of three
triples, which go to Phase~3, or write \(0\) on tape~\(2\) and move the heads
of tapes \(2\) right and \(3\) left, or do the same with \(1\).  Each written
bit moves the head of tape~\(3\) one cell nearer to cell \(0\), so at most
\(n\) bits are written, and every bit string \(y\) with
\(\operatorname{len}(y)\leq n\) is written on some branch.  Tape~\(2\) then
holds \(\langle x,y\rangle\).  This takes at most \(n+2\) steps.

\emph{Phase 3.}  It returns the head of tape~\(2\) to cell \(0\), in at most
\(3n+2\) steps, since \(\operatorname{len}(\langle x,y\rangle)\leq 3n+2\).

\emph{Phase 4.}  It simulates \(M_R\) one step per step, with tape~\(2\) in
the role of tape~\(1\) of \(M_R\) and tapes \(4,\ldots,\kappa+2\) in the
roles of tapes \(2,\ldots,\kappa\); these tapes are blank with heads on cell
\(0\), and cell \(0\) of tape~\(2\) is read unmarked.  It accepts or rejects
as \(M_R\) does on \(\langle x,y\rangle\), within
\(p(\operatorname{len}(\langle x,y\rangle))\leq p(3n+2)\) steps.

Phases~1 and~3 are deterministic and Phase~4 is deterministic given \(y\), so
the branches of \(N\) on \(x\) correspond to the choices in Phase~2, that is,
to the strings \(y\) with \(\operatorname{len}(y)\leq n\).  Every branch
halts within \(g(n)=(2n+3)+(n+2)+(3n+2)+p(3n+2)\) steps.  The branch for
\(y\) accepts iff \(R(x,y)\).  By hypothesis \(x\in L\) iff some \(y\) with
\(\operatorname{len}(y)\leq n\) has \(R(x,y)\), so \(N\) accepts exactly the
members of \(L\), and every branch halts, so \(N\) decides \(L\).  For
\(n\geq1\), \(g(n)\leq 6n+7+e\,(3n+3)^d\leq(13+e\,6^d)\,n^d\), so \(L\in
\mathsf{NTIME}(n^d)\subseteq\mathsf{NP}\) in the sense of
Definition~\ref{def:p04-resource-classes}.
\end{proof}
\end{annexitem}

\begin{annexitem}{annex:p02-w02b-p02-015-cook-levin}{thm:p04-cook-levin}{Cook--Levin theorem}
\begin{proof}
The language \(\mathsf{SAT}\) of the theorem is that of
Definition~\ref{def:p02-sat-taut} (\cref{sec:p04-complexity-complete-problems}),
with the code table fixed in Convention~\ref{def:p02-w02-p02-018-input}.
Reductions and classes are those of
Convention~\ref{def:p04-complexity-convention} and
Definition~\ref{def:p04-resource-classes}, with machines read as in
Convention~\ref{def:p02-w02b-p02-015-machine-reading}.

\emph{Membership.}  Let \(R\) be the predicate of the proof of
Theorem~\ref{thm:p02-w02-p02-018-membership}: \(R(x,y)\) holds iff \(x\)
encodes a formula \(\varphi\) with \(k\) distinct variables, listed in order
of first occurrence, \(\operatorname{len}(y)=k\), and the valuation read off
\(y\) satisfies \(\varphi\).  By that theorem and its proof in
Annex~\ref{annex:p02-w02-p02-018-membership}, \(R\) is decided on
\(\langle x,y\rangle\) in a number of steps polynomial in the length of
\(\langle x,y\rangle\), and \(x\in\mathsf{SAT}\) iff some \(y\) satisfies
\(R(x,y)\).  Any such \(y\) has length \(k\), and \(k\leq n\), where \(n\) is the
bit length of \(x\), by the bound on the number of distinct variables in
Convention~\ref{def:p02-w02-p02-018-input}.  So \(x\in\mathsf{SAT}\) iff some \(y\) with
\(\operatorname{len}(y)\leq n\) satisfies \(R(x,y)\), and clause~3 of
Lemma~\ref{thm:p02-w02b-p02-015-bridge} gives
\(\mathsf{SAT}\in\mathsf{NP}\) in the sense of
Definition~\ref{def:p04-resource-classes}.

\emph{Hardness.}  Let \(L\in\mathsf{NP}\) in the sense of
Definition~\ref{def:p04-resource-classes}.  By clause~2 of
Lemma~\ref{thm:p02-w02b-p02-015-bridge}, \(L\in\mathsf{NP}\) in the sense of
Convention~\ref{def:p02-complexity-standing}, with a predicate as in
Convention~\ref{def:p02-w02-p02-018-verifier}.  By
Theorem~\ref{thm:p02-w02-p02-018-cook-levin} there is a map
\(x\mapsto\varphi_x\), computable in polynomial time by a deterministic
multitape machine, with \(\varphi_x\) an encoded formula and \(x\in L\) iff
\(\varphi_x\) is satisfiable, that is, iff the encoding of \(\varphi_x\)
belongs to \(\mathsf{SAT}\).  By clause~1 of the lemma the same map is
computed in polynomial time by a deterministic machine with read-only input,
so it witnesses \(L\leq_m^p\mathsf{SAT}\) in the sense of
Convention~\ref{def:p04-complexity-convention}.

\emph{Completeness.}  \(\mathsf{SAT}\) is a member of \(\mathsf{NP}\), and
every language of \(\mathsf{NP}\) reduces to it by \(\leq_m^p\); by
Convention~\ref{def:p04-complexity-convention} this is
\(\mathsf{NP}\)-completeness.
\end{proof}
\end{annexitem}
